\documentclass[tikz,border=6pt]{standalone}
\usepackage{ctex}
\usepackage{amsmath}
\usetikzlibrary{calc, arrows.meta}

\begin{document}
\begin{tikzpicture}[scale=1.0, thick]

% 相关点法示例：M 在圆 x^2+y^2=4 上运动，P 是 OM 的中点
% M(x0,y0) 满足 x0^2+y0^2=4
% P 是 OM 中点：P=(x0/2, y0/2)，设 P=(x,y)
% 则 x0=2x, y0=2y，代入：(2x)^2+(2y)^2=4 → x^2+y^2=1（圆）

% 坐标轴
\draw[->, gray, thin] (-2.8, 0) -- (2.8, 0) node[right, font=\footnotesize] {$x$};
\draw[->, gray, thin] (0, -2.8) -- (0, 2.8) node[above, font=\footnotesize] {$y$};
\node[font=\footnotesize, below left] at (0,0) {$O$};

% 原曲线：圆 x^2+y^2=4（M 在此运动）
\draw[blue, thick] (0,0) circle [radius=2.0];
\node[font=\footnotesize, blue] at (1.5, 1.75) {$M$ 所在曲线};
\node[font=\footnotesize, blue] at (1.5, 1.4) {$x_M^2+y_M^2=4$};

% 轨迹：圆 x^2+y^2=1（P 的轨迹）
\draw[red, thick] (0,0) circle [radius=1.0];
\node[font=\footnotesize, red] at (0.7, -1.05) {$P$ 的轨迹};
\node[font=\footnotesize, red] at (0.7, -1.38) {$x^2+y^2=1$};

% 示例点 M1（角度 50°）
\coordinate (M1) at ({2*cos(50)}, {2*sin(50)});
\coordinate (P1) at ({cos(50)}, {sin(50)});

% 连线 O-M1
\draw[dashed, gray] (0,0) -- (M1);

% 标注 M1, P1
\fill[blue] (M1) circle [radius=2pt];
\fill[red] (P1) circle [radius=2pt];
\node[font=\footnotesize, blue, above right] at (M1) {$M(x_M,y_M)$};
\node[font=\footnotesize, red, below right] at (P1) {$P(x,y)$（中点）};

% 示例点 M2（角度 200°）
\coordinate (M2) at ({2*cos(200)}, {2*sin(200)});
\coordinate (P2) at ({cos(200)}, {sin(200)});
\draw[dashed, gray] (0,0) -- (M2);
\fill[blue] (M2) circle [radius=2pt];
\fill[red] (P2) circle [radius=2pt];
\node[font=\footnotesize, blue, left] at (M2) {$M'$};
\node[font=\footnotesize, red, right] at (P2) {$P'$};

% 相关点法步骤标注（右下）
\node[draw, rounded corners, fill=yellow!10, font=\footnotesize, align=left]
  at (2.2, -1.9)
  {相关点法步骤：\\
   1. 设 $P(x,y)$，$M$ 与 $P$ 关系\\
   2. 用 $x,y$ 表示 $x_M,y_M$\\
   3. 代入 $M$ 的曲线方程\\
   4. 化简得 $P$ 的轨迹};

% 中点关系标注
\node[font=\footnotesize, gray] at (0.65, 0.9) {$OP{=}\tfrac{1}{2}OM$};

\end{tikzpicture}
\end{document}
